\documentclass[a4paper]{article}

%% Useful packages
\usepackage{amsmath}
\usepackage{graphicx}	
\usepackage{algpseudocode}
\usepackage{algorithm}
\usepackage[colorinlistoftodos]{todonotes}
\usepackage[colorlinks=true, allcolors=black]{hyperref}
\usepackage[fontsize=11pt]{scrextend}
\usepackage{titlesec}
 \usepackage[section]{placeins}

\setlength\parindent{0pt}
\titleformat*{\section}{\Large\bfseries}
\titleformat*{\subsection}{\Large\bfseries}

\title{PowerEnjoy Service - Integration Test Plan Document}

\begin{document}

\begin{titlepage}
\begin{figure}
\centering
\includegraphics[width=0.2\textwidth]{polimi.jpg}
\par
\LARGE Politecnico di Milano
\end{figure}


\maketitle
\textbf{Version 1.1}
\newline

\raggedright
Authors:
\begin{itemize}
	\item Domenico FAVARO (Mat. 837995)
        	\item Matheus FIM (Mat. 876069)
	\item Caio ZULIANI (Mat. 877266)	
\end{itemize}
\raggedleft
Prof. Elisabetta DI NITTO
\thispagestyle{empty}
\end{titlepage}

\tableofcontents
\newpage
 
\section{Introduction}
\subsection{Revision History}
This section records all revisions to the Document.
\newline \newline
\begin{tabular}{ | c | c | c | c | }
\hline
	Version & Date & Authors & Summary \\ \hline
	1.1 & 15/01/16 & Domenico Favaro, Caio Zuliani, Matheus Fim & Initial Release  \\ \hline
\end{tabular}

\subsection{Purpose and Scope}
The Integration Test Plan Document (ITPD) serves to present the integration sequence and testing for all subsystems and components that conform PowerEnjoy Car Sharing Service.

\subsection{Definitions and Abbreviations}
\begin{itemize}
\item \textbf{RASD:} Reqirements And Specifications Document.
\item \textbf{DD:} Design Document.
\item \textbf{ITPD:} Integration Test Plan Document.
\item \textbf{ETC..:}
\end{itemize}
For other concepts concerning the Service definition look in the \textbf{Glossary} section of the RASD and DD.

\subsection{Reference Documents}
\begin{itemize}
\item Specification Document: Assignments AA 2016-2017.pdf
\item PowerEnjoy Requirements And Specifications Document (RASD)
\item PowerEnjoy Design Document (DD)
\item Example Document - Integration testing example document.pdf
\item Testing Tools Documents:
\begin{itemize}
\item[-] Mockito
\item[-] JMeter
\end{itemize}
\end{itemize}

\newpage
\section{Integration Criteria}
\subsection{Entry Criteria}
Specify the criteria that must be met before integration testing of specific elements may begin (e.g., functions must have been unit tested).

\subsection{Elements to be Integrated}
Identify the components to be integrated, refer to your design document to identify such components in a way that is consistent with your design.

\subsection{Integration Testing Strategy}
Describe the integration testing approach (top-down, bottom-up, functional groupings, etc.) and the rationale for the choosing that approach.

\subsection{Sequence of Component/Function Integration}
\subsubsection{Software Integration Sequence}
For each subsystem, identify the sequence in which the software components will be integrated within the subsystem; relate this sequence to any product features that are being build up.
\subsubsection{Subsystem Integration Sequence}
Identify the order in which subsystems will be integrated.

\newpage
\section{Individual Steps and Test Description}

\section{Performance Analysis}

\section{Required Tools and Test Equpment}
\subsection{Tools}
\subsection{Test Equipment}

\section{Required Program Stubs and Test Data}
\subsection{Program Stubs}
\subsection{Test Data}

\newpage
\section{Effort Spent}
\begin{tabular}{ | c | c | c | c | }
\hline
	\textbf {Date} & \textbf {Domenico} & \textbf {Caio} & \textbf {Matheus} \\ \hline
	27/12/16& 2h & 2h & 2h  \\ \hline
	28/12/16& - & - & - \\ \hline
	29/12/16& 1h & - & - \\ \hline
	30/12/16& 2h & - & - \\ \hline
	31/12/16& - & - & - \\ \hline
	01/01/17& - & - & - \\ \hline
	02/01/17& - & - & - \\ \hline
	03/01/17& - & - & - \\ \hline
	04/01/17& - & - & - \\ \hline
	05/01/17& - & - & - \\ \hline
	06/01/17& - & - & - \\ \hline
	07/01/17& - & - & - \\ \hline
	08/01/17& - & - & - \\ \hline
	09/01/17& - & - & - \\ \hline
	10/01/17& - & - & - \\ \hline
	11/01/17& - & - & - \\ \hline
	12/01/17& - & - & - \\ \hline
	13/01/17& - & - & - \\ \hline
	14/01/17& - & - & - \\ \hline
\end{tabular}
\newpage

\section{Changelog}
As the project and design decisions may change during the development this document is also prone to change.
We'll document every version in this part.
\begin{itemize}
\item \textbf {Version 1.1:} 15/01/2017
\end{itemize}
\end{document}
